% !PRINTED 560
\weberpage{560}
%<<CONT
from the remark that the substitution $(r, r_h)$ calls forth among the $\eta_k$
the substitution $(\eta_k, \eta_{k+h})$, when the index of $\eta$
is taken modulo $e$. The group of $F_e(x) = 0$
therefore consists of the powers of the cyclic permutation

\[
(\eta, \eta_1, \eta_2 \ldots \eta_{e-1}).
\]

\begin{center}
\S.\,168.

\gesperrt{The Gauss Method for the Computation of the\\
Resolvents.}
\end{center}

In order actually to carry out the solution of the cyclotomic
equation in concrete cases, all that matters, after what we
have developed in the preceding section, is to represent the
products of any two conjugate periods as linear functions
of these same periods. For this \gesperrt{Gauss} has given a simple
procedure, which we shall now become acquainted with. To this
end we must alter the notation of the periods a little.

We write any two of the periods \S.\,167, (12) in
the form

\weq{1}{\begin{aligned}
\eta^{(\lambda)} &= r^\lambda + r^{\lambda'} + r^{\lambda''} + \cdots\\
\eta^{(\mu)} &= r^\mu + r^{\mu'} + r^{\mu''} + \cdots
\end{aligned}}

so that, if $\lambda \equiv g^h \ (\mathrm{mod}\ n)$, then $\eta^{(\lambda)} = \eta_h$ or also
$\eta^{(\lambda)} = \eta_{\mathrm{ind}\,\lambda}$; we then have

\weq{2}{\begin{aligned}
\lambda' &\equiv \lambda g^e, & \lambda'' &\equiv \lambda g^{2e} \cdots\\
\mu' &\equiv \mu g^e, & \mu'' &\equiv \mu g^{2e} \cdots
\end{aligned}\ (\mathrm{mod}\ n);}

if, then, we let the symbols $s$ and $t$ each run through a complete system of residues
modulo $f$, we have

\[
\eta^{(\lambda)} = \sum^{s} r^{\lambda g^{se}}, \qquad \eta^{(\mu)} = \sum^{t} r^{\mu g^{te}}.
\]

The product of these is

\[
\eta^{(\lambda)} \eta^{(\mu)} = \sum^{s}\sum^{t} r^{\lambda g^{se} + \mu g^{te}}.
\]

If, in forming this sum, we first hold $s$ fixed,
and sum with respect to $t$, we may replace $t$ by $t + s$, because both run simultaneously through a complete system of residues
modulo $f$. Accordingly we obtain
\[
\eta^{(\lambda)}\eta^{(\mu)} = \sum^{s}\sum^{t} r^{(\lambda + \mu g^{te})g^{se}}.
\]
%CONT>>
